\section{A longitudinal observatory for machine research}

\subsection{The live mathematical case study}

The observational substrate is a live research programme in \texttt{SzeChunYiu/RAKL\_math}. The repository is deliberately separated from the reusable \RAKL{} framework repository. The application repository contains problem contracts, context fibres, research DAGs, source packets, candidate artifacts, falsifiers, reviews, traces and problem-specific tests for six unsolved Millennium Prize Problems: P versus NP, the Riemann Hypothesis, Navier--Stokes existence and smoothness, Yang--Mills existence and mass gap, the Hodge Conjecture, and the Birch and Swinnerton--Dyer Conjecture. A cross-problem lane studies transfer and recurring process failures across these domains.

At the initial Paper 5 freeze, the active scheduled programme contains nine hourly research lanes: P versus NP; an analytic Riemann-Hypothesis lane; Navier--Stokes regularity and blow-up lanes; constructive and spectral/RG Yang--Mills lanes; a Hodge deformation lane; an analytic BSD lane; and a cross-Millennium transfer/meta-observer lane. Complementary spectral RH, motive-oriented Hodge and Selmer/Iwasawa BSD lanes are preserved but paused. The schedule is an operational fact, not a claim of nine independent replicates. All lanes share repository state and can influence one another through persistent memory, so statistical analyses must treat them as coupled longitudinal processes.

Every application run is instructed to read current \RAKL{} \texttt{main} as the framework source of truth, while new mathematical application work remains in \texttt{RAKL\_math}. This split serves two purposes. First, problem evidence does not silently become framework authority. Second, a framework defect exposed by one domain can be opened as a separate \RAKL{} issue and evaluated as a method hypothesis rather than patched inside the mathematical result that exposed it.

\subsection{Observation is not hidden reasoning disclosure}

The observatory does not attempt to reconstruct private model chain-of-thought. The recorded object is an auditable scientific decision trace: problem state, retrieved evidence, proposed action, executed tools, observable result, verification, residual and next action. This distinction is methodologically important. A reproducible research record need not contain the inaccessible or potentially unstable private reasoning process that generated a proposal.

\begin{definition}[Task episode]
A task episode is an immutable externally observable record
\[
E=(p,a,c,F,O,\tau,V,y,r,\kappa,t,h),
\]
where $p$ is the problem/atom identity, $a$ the active atomic target, $c$ the context identity, $F$ the problem-fibre snapshot, $O$ the operators or methods attempted, $\tau$ the public action/observation trace, $V$ the verification evidence, $y$ the outcome, $r$ the residual signature, $\kappa$ a resource/cost record, $t$ the timestamp and $h$ the content identity.
\end{definition}

A non-success outcome should carry an explicit residual or blocker. The raw episode is never replaced by a later reflection summary.

\subsection{Episode, diagnosis and lesson are different objects}

A central design decision is
\[
\boxed{\text{episode} \neq \text{diagnosis} \neq \text{obstruction} \neq \text{lesson}}.
\]
An episode records what happened. A diagnosis proposes why it happened. An obstruction generalizes a supported boundary that may apply elsewhere. A lesson proposes a reusable action under explicit triggers and scope conditions. Conflating these stages causes a familiar failure in autonomous agents: a fluent post-hoc explanation is treated as a verified causal rule.

A versioned lesson therefore records
\[
L=(g,c,a,e,b,S^+,S^-,f,\alpha,\nu,h_L),
\]
with trigger $g$, scope $c$, proposed action $a$, expected effects $e$, boundaries $b$, supporting and contradicting episodes $S^+,S^-$, a falsifier $f$, authority $\alpha$, validation obligations $\nu$ and content identity $h_L$. Reusable promotion requires evidence beyond reflection. Raw episodes remain evidence roots even when a lesson becomes useful.

\subsection{Problem-conditioned fibres}

A research agent should not retrieve the entire archive for every action. For active problem atom $A$ under context $c$, \RAKL{} compiles a problem-conditioned fibre
\[
\mathcal F(A\mid P,c)=
(K_A,O_A,E_A,B_A,M_A,X_A),
\]
where $K_A$ contains epistemic items, $O_A$ applicable tools/operators, $E_A$ analogous episodes, $B_A$ failures and boundary warnings, $M_A$ strategy motifs, and $X_A$ expertise or unresolved warnings. Co-retrieval does not imply compatibility or authority.

The fibre snapshot is important for retrospective method analysis. If a relevant item existed in the bound search universe but was not retrieved, the failure is different from a case in which the item was not present. If the item was retrieved but rejected, the rejection rationale becomes part of the observable decision record.

\subsection{Local results and gluing failures}

Many difficult research episodes do not fail inside a local lemma. They fail because locally valid pieces cannot be assembled into a root result. \RAKL{} represents this explicitly. A decomposition produces problem atoms with dependencies and optional interface keys. Local sections can be verified independently, but a global solution requires complete coverage, dependency closure and compatibility on shared interfaces.

We therefore distinguish
\[
\text{local mathematical failure}
\]
from
\[
\text{gluing/interface failure}.
\]
The latter includes missing uniformity across scales, a theorem whose hypotheses are not inherited by the target limit, a representation that does not preserve a root-critical coordinate, or a proof component that is true only in a different category. Failed gluing is itself an experience record and can become a recurring obstruction family.

\subsection{Vector saturation}

Repeated unsuccessful search should not be summarized as ``stuck.'' The v3 saturation state tracks bounded novelty across seven coordinates:
\[
\operatorname{Sat}(R,D)=
(S_K,S_O,S_E,S_B,S_R,S_P,S_M),
\]
corresponding to knowledge, operator, experience-pattern, obstruction, relation, path and meta-method novelty. An axis is considered bounded-flat only under a registered search window with sufficient independent route coverage, no retained novelty on that axis and no native residual that reopens it. Saturation never licenses an absolute completeness claim.

The observatory records which axis was reopened after a consequential episode. This permits later questions such as whether an agent wastes cycles by continuing to search the knowledge axis when the stable residual is representational or relational.

\subsection{Structural novelty of a solved subproblem}

A verified local result can be classified by the strongest genuinely new problem-solving structure required:

\begin{center}
\begin{tabular}{ll}
\toprule
Class & Interpretation \\
\midrule
\texttt{STORED} & exact verified solution already registered \\
\texttt{RAKL\_TRIVIAL} & composition of existing resources/operators \\
\texttt{TRANSFER\_NOVEL} & new application of existing structure with mapping witness \\
\texttt{REPRESENTATION\_NOVEL} & new representation/bridge/invariant required \\
\texttt{OPERATOR\_NOVEL} & genuinely new transformation/operator required \\
\texttt{ONTOLOGY\_NOVEL} & schema/ontology extension required \\
\texttt{UNRESOLVED} & verification or ancestry insufficient \\
\bottomrule
\end{tabular}
\end{center}

This classification supports an empirical question that cannot be answered from final problem labels alone: what fraction of useful research progress requires genuinely new problem-solving primitives, and what fraction is retrieval, transfer, recomposition, representation repair or gluing?

\subsection{Method-case telemetry}

For Paper 5, each new consequential research cycle is instructed to emit a bounded \texttt{RAKL\_METHOD\_CASE\_STUDY} record. The record contains:

\begin{itemize}[leftmargin=*]
\item the research method or operator sequence actually used;
\item which prior successes and failures affected routing;
\item relevant items retrieved but rejected, and relevant items later discovered to have been missed;
\item the chosen falsifier or verification action;
\item outcome and residual;
\item whether the failure is primarily mathematical, retrieval, decomposition, representation, gluing, source/provenance, verification, tooling/CI or meta-policy;
\item the saturation coordinate reopened;
\item the strongest defensible novelty class;
\item a framework-improvement hypothesis, if one is warranted, with a cheapest prospective discriminator.
\end{itemize}

The method-case record is an observational layer. It does not authorize a framework change.

\subsection{Human and infrastructure interventions}

The longitudinal study also records exogenous interventions. Examples include manual schedule-prompt changes, human merge overrides, model/tool upgrades, CI changes and repository migrations. Without an intervention ledger, a later performance change can be falsely attributed to experience learning when it actually came from a stronger base model or a changed tool surface.

Accordingly, a version comparison should bind at least model identity/revision, system-prompt or harness identity, tool policy, resource ceilings, repository state and evaluator identity. If a provider changes an unpinnable model behind an endpoint, comparisons spanning that change are treated as a new observational epoch unless a paired same-time control can recover comparability.
